\subsection{The secondary root approaching the central line}
\label{sec:direct-negative-weak-root}

The remaining part of the coefficient cover has \(a\downarrow0\).
Normalization in \eqref{eq:direct-negative-secondary-scales} then
forces \(b=1\). To treat the positive root approaching the central
line, apply \cref{prop:weak-root-passage} with
\[
 \kappa=s/a,\qquad \ell=sra,\qquad \hat r=r,\qquad \epsilon_2=a.
\]
These choices give \eqref{eq:weak-root-scales} and the logarithmic
height coordinate \(T=-s\log h\). The remaining hypothesis concerns
the negative branch outside the local weak-root region. We verify it
by separating bounded secondary outputs from graphs that reach the
annulus while still thin. The first group has a positive endpoint
rate; for the second, the annular and far-root estimates give the
uniform lower bound required by the weak-root comparison.

\begin{proposition}[Negative quadratic drift at the weak root]
\label{prop:direct-negative-weak-root}
There are fixed \(a_-,\kappa_->0\) such that \(\Delta_p\) has at
most two zeros on
\[
 0<a\le a_-,\qquad 0<s/a\le\kappa_-
\]
for all sufficiently small independent \(s,r>0\). Consequently, for all
sufficiently small independent \(s,r>0\) and every \(L>1/\kappa_-\),
\(\Delta_p\) has at most three zeros on the part of
\eqref{eq:direct-negative-secondary-scales} outside the strip
\(|q_0|\le Lq_1^2\).
\end{proposition}

\begin{proof}
At the radius \(\ell_d=srd\), the outward drift on the negative branch is
\[
 K_d(x)=a_d+x+dx^2+rd^2x^3R_q(-\ell_dx).
\]
For a fixed large \(M\) and \(r\le r(M)\),
\eqref{eq:weak-root-negative-drift} holds with \(d_{\max}=M\), with
constants depending on \(M\). We verify
\eqref{eq:weak-root-outer} for the negative graphs with
\(z<Y_b|\eta|^2\) at \(|\eta|=srM\).

Let \(X=|\eta|/(sr)\), and follow such a graph to the height
\(\ell_M^2Y_b\). At \(d=M\), the local radius in the weak-root estimate
is \(x=X/M\). If \(X\le4M\) at this height, the corresponding projected
state lies in \([0.9,4]\), and
\cref{lem:weak-root-branches}(ii) with \(d=M\) gives
\(\mathfrak H_-\ge c/M\) and \(|I_{-,T}|=o(1)\).

If \(X>4M\) at that height, the graph is thin at the annular cut \(X=M\):
\[
 z(-srM;h)/(s^3r^2M^2)<\delta_{\rm an},
\]
where \(\delta_{\rm an}\) is the thin threshold of
\cref{sec:direct-negative-annulus}. To obtain this estimate, use
\(|\eta|=\ell_Mx\), \(z=\ell_M^2Y\) with the same \(d=M\).
The outward graph reaches \(x=4\) below \(Y_b\) and lies in the box between
\(x=1\) and \(x=4\). Also, \(8B_0Y_b\le15\) since \(Y_b\) is
small. For small \(s\), apply \cref{lem:direct-radial-advance} in
\(x\), with \(\delta=s/M\), \(X_1=1\), \(X_2=4\) and
\(C_F=\sup|K_M/x|\) on the box. The height at \(x=1\) is at most
\(Y_b\e^{-3M/(2C_Fs)}<s\delta_{\rm an}\), as required.

For this thin input, take \(M\) large enough for the annular margins
on the negative branch. In the secondary coordinate \(X=|\eta|/(sr)\),
the negative drift is comparable to \(a+X\) and its derivative is
bounded below on this fixed box. Here
\(\delta=s\) and the weak-root condition gives \(a/s\ge1/\kappa_-\);
\cref{lem:direct-source-height} supplies the height unit at a fixed
secondary cut below \(M\), uniformly as \(a\to0\).
\Cref{lem:negative-branch-far-root}, whose constants are uniform down
to \(a=0\), gives, for exits of the negative branch up to \(\rho=0.9\),
\[
 \mathfrak H_-=2-3\rho+
                O(r/\rho+s(1+|\log s|))\ge-C_n,\qquad
 |I_{-,T}|\le C(s\rho)^{2-4s}=o(1).
\]
Graphs still thin at \(\rho=0.9\) satisfy \(\mathfrak H_-\ge-C_n\) and
\(|I_{-,T}|\le Cs^2\). Thus \eqref{eq:weak-root-outer} holds in both
cases. \Cref{prop:weak-root-passage} gives at most two zeros of
\(\Delta_p\), with \(a_-\) and \(\kappa_-\) equal to its constants
\(a_\circ\) and \(\kappa_\circ\).

To obtain the assertion outside the contact strip, combine the weak-root
bound with the compact coefficient cover of
\cref{sec:direct-negative-separated}, which gives the bound three on
\(a\ge a_-/2\), overlapping the weak-root region. On the remaining
small-\(a\) part, where \(s/a>\kappa_-\), normalization gives
\(b=1\) and \(|q_0|/q_1^2=a/s<1/\kappa_-\). This part lies in
\(|q_0|\le Lq_1^2\) for every \(L>1/\kappa_-\).
Hence the weak-root and compact-coefficient regions cover the parameter
set outside the contact strip. Their bounds on the whole interval
\(\mathcal H_p\) give at most three zeros there.
The thresholds for \(s,r\) are those of these two regions and are
independent of \(L\).
\end{proof}
